\documentclass[11pt]{article}
\usepackage[a4paper,margin=1in]{geometry}
\usepackage[T1]{fontenc}
\usepackage{lmodern}
\usepackage{amsmath,amssymb,amsthm}
\usepackage{booktabs}
\usepackage{array}
\usepackage[hidelinks]{hyperref}
\usepackage{microtype}

\newtheorem{theorem}{Theorem}
\newtheorem{lemma}[theorem]{Lemma}
\newtheorem{proposition}[theorem]{Proposition}
\newtheorem{corollary}[theorem]{Corollary}
\theoremstyle{remark}
\newtheorem{remark}[theorem]{Remark}

\newcommand{\st}[1]{\mathrm{#1}}
\newcommand{\w}[1]{\mathtt{#1}}
\newcommand{\inftyl}{0^{\infty}}
\newcommand{\coq}[1]{\texttt{#1}}
\newcommand{\up}{\uparrow}

\title{Exact Step Counts and an Exact Halting Time\\ for the BB(6) Champion, Proved in Coq}
\author{Vishnu Vaidyanathan\\
\small \texttt{vishnu@edgi.in}\\
\small ORCID 0009-0001-6244-5262\\
\small \url{https://github.com/VishnuVaidyanathan/BB5-decimal-macrogrammar}}
\date{November 2026}

\begin{document}
\maketitle

\begin{abstract}
The current 6-state Busy Beaver champion, found by mxdys in 2025, runs for more
than $2\up\up\up 5$ steps. mxdys proved in Coq that it halts, using a small set
of rules of the form ``this tape shape reaches that tape shape''. Those rules
carry no step counts, and the halting time is known only as an estimate of its
size. We add the missing step counts. For every move of the machine (the tick
of its binary counter, a full countdown, the start of an epoch, the restart,
and the final move) we give the exact number of steps as a short formula, valid
for all values of the parameters. From these we build an exact expression for
the halting time and for the number of ones on the final tape. We prove in Coq,
with no axioms, that the machine started on the blank tape halts after exactly
that many steps and not before, and that the final tape holds exactly that many
ones. The numbers cannot be written out, but the theorems fix which numbers they
are. The halting time is close to two thirds of the square of the number of ones. The moves were first found by reading the transcript of the machine, the
method of our earlier paper on the BB(5) champion.
\end{abstract}

\section{Introduction}

The machine studied here is
\[
\texttt{1RB1RA\_1RC1RZ\_1LD0RF\_1RA0LE\_0LD1RC\_1RA0RE},
\]
found by mxdys in June 2025. It is the current BB(6) champion: it halts, and
its running time and its number of ones lie between $2\up\up 2\up\up 2\up\up 10$
and $2\up\up 2\up\up 2\up\up 11$, which is more than $2\up\up\up 5$
\cite{wiki}. The proof that it halts is a Coq development by mxdys \cite{mxdys}.
That proof describes the tape by a few numbers and shows how these numbers
change. Every statement in it has the form $c \to^{*} c'$: the configuration
$c$ reaches $c'$ after \emph{some} number of steps. The number of steps is never
stated, and so the proof does not say when the machine halts.

This paper supplies the step counts.

\paragraph{What is not new.} The tape shapes, the rules that take one shape to
the next, the fact that the machine counts in binary, the size of the running
time and the proof of halting are all the work of mxdys \cite{mxdys,wiki}. We
found the same moves independently by reading the transcript of the machine,
but they are the same moves.

\paragraph{Contributions.}
\begin{enumerate}
\item An exact step count for each move of the machine, for all values of its
parameters (Theorems \ref{thm:tick}--\ref{thm:halt}). Each is a linear formula,
except the full countdown, which is a quadratic in $2^{b}$.
\item An exact expression for the halting time, and a Coq proof that the
machine halts after exactly that many steps and at no earlier step
(Theorem~\ref{thm:main}).
\item An exact expression for the number of ones on the final tape, also proved
in Coq (Theorem~\ref{thm:ones}).
\item The transcript words of each move (Section~\ref{sec:words}). These are
tested against simulation but are not yet stated in Coq.
\end{enumerate}
The Coq file is self-contained (about 890 lines, Coq 8.18, no axioms). It does
not import the development of mxdys; it proves the moves again in its own small
model, now with step counts.

\section{Preliminaries}

\paragraph{The machine.} The transition table is

\begin{center}
\begin{tabular}{lcccccc}
\toprule
 & A & B & C & D & E & F\\
\midrule
read 0 & 1RB & 1RC & 1LD & 1RA & 0LD & 1RA\\
read 1 & 1RA & 1RZ & 0RF & 0LE & 1RC & 0RE\\
\bottomrule
\end{tabular}
\end{center}

\noindent The machine starts in state A on a blank tape. The only halting pair
is B reading 1. We count the halting transition as a step, as in \cite{bb5}.

\paragraph{Configurations.} We write $\inftyl\, u\, [\st{s}\ x]\, v\, \inftyl$ for
the tape $u\,x\,v$ surrounded by blanks, with the head in state $\st{s}$ on the
cell holding $x$. Powers denote repetition.

\paragraph{Transcripts.} The transcript of a run is the list of state--symbol
pairs it visits \cite{bbchallenge}. The machine has twelve pairs, so we write
each as one character,
$d(\st{s},x)=2\cdot\mathrm{idx}(\st{s})+x$, with the values ten and eleven written
$\w{a}$ and $\w{b}$:
$\st{A}0=\w0$, $\st{A}1=\w1$, \dots, $\st{E}1=\w9$, $\st{F}0=\w{a}$, $\st{F}1=\w{b}$.
The halting pair $\st{B}1$ is the digit $\w3$.

\paragraph{Three basic words.} Almost every step of the run belongs to one of
three repeated words.
\begin{itemize}
\item $\w1$ ($\st{A}1$): the head walks right over a one.
\item $\w{479}$ ($\st{C}0,\st{D}1,\st{E}1$): a $0$ under the head moves one cell
to the left through a block of ones:
$\ 1\,1\,[\st{C}\ 0] \to 1\,[\st{C}\ 0]\,1$.
\item $\w{5b9}$ ($\st{C}1,\st{F}1,\st{E}1$): the block $111$ becomes $001$ and the
head moves three cells to the right.
\end{itemize}

\section{Tape shapes}\label{sec:shapes}

The shapes below are those of mxdys \cite{mxdys}, written for this machine.

\paragraph{Counters.} A counter is a number written in binary, high bit first,
with the bit $1$ written as the block $100$ and the bit $0$ written as the
block $111$. Counters count down.

\paragraph{The inner shape $S_0$.} A counter with any tape $X$ to its left, the
head in state E on the last cell of the counter, and to the right the block
\[
R(b) \;=\; 0\ 11100\ (111)^{1+b}\ \inftyl .
\]

\paragraph{The outer shape $S_1$.} The whole tape at the start of an epoch is
\[
S_1(\beta, q, b) \;=\; \inftyl\ 1110\,11\ \ \beta\ \ 100\ \ 10\ \ 1^{q-1}\,[\st{E}\ 1]\ \ R(b),
\]
where $\beta$ is a counter, called the \emph{outer counter}, and $q$ is the
number of ones after the marker $10$. (For $q=0$ the head is on the $0$ of the
marker.) In the numbers of mxdys, $S_1(\mathit{len}_0,a_0,m,a,0,b)$, the outer
counter has $\mathit{len}_0$ bits and value $a_0$, and
\[
q = m + 3a, \qquad m\in\{0,1,2\}.
\]
So one number $q$ carries both the marker $m = q \bmod 3$ and the length $a$ of
the inner counter, which is all zeros at the start of an epoch. This makes the
rules shorter.

Throughout, $N_k = 2^{k}-1$.

\section{The tick}

\begin{theorem}[Tick; \coq{tick}]\label{thm:tick}
Let a counter end in a bit $1$ followed by $j \ge 0$ bits $0$. From the shape
$S_0$ with this counter and right block $R(b)$, the machine reaches the shape
$S_0$ in which that bit $1$ has become $0$, the $j$ bits after it have become
$1$, and the right block is $R(b+1)$. It takes exactly
\[
12b + 12j + 48 \quad \text{steps.}
\]
\end{theorem}

In other words the counter goes down by one, the right block gains one $111$,
and $j$ is the length of the borrow. This is the lemma \coq{Inc0} of
\cite{mxdys}, with the step count added.

\begin{theorem}[Countdown; \coq{full\_ticks}, \coq{T\_closed}]\label{thm:count}
Let the lowest $k$ bits of a counter all be $1$. From $S_0$ with right block
$R(b)$ the machine reaches $S_0$ with those $k$ bits all $0$ and right block
$R(b+N_k)$, in exactly $T(k,b)$ steps, where
\[
T(k,b) \;=\; 6N_k^{2} + 54N_k + 12\,b\,N_k - 12k .
\]
\end{theorem}

\begin{proof}[Idea of the proof]
To clear $k+1$ bits, clear the low $k$ bits, do one tick with borrow $k$, and
clear the low $k$ bits again. This gives the recursion
$T(k+1,b) = T(k,b) + \bigl(12(b+N_k)+12k+48\bigr) + T(k,\,b+N_k+1)$ with
$T(0,b)=0$, and the closed form follows by induction on $k$. No sum over the
borrow lengths is needed.
\end{proof}

\section{The epoch}

An epoch starts at a shape $S_1$ whose outer counter is not zero.

\begin{theorem}[Epoch; \coq{load}, \coq{epoch}]\label{thm:epoch}
Let the outer counter $\beta$ end in a bit $1$ followed by $j_0\ge 0$ bits $0$,
and let $\beta'$ be $\beta$ minus one. Then
\[
S_1(\beta, q, b) \;\longrightarrow\; S_1\bigl(\beta',\ q+3b+5,\ 2^{b}-1\bigr)
\]
in exactly
\[
\underbrace{4q + 12j_0 + 3b + 63}_{\text{load}} \;+\; \underbrace{T(b,0)}_{2^{b}-1\ \text{ticks}}
\quad\text{steps, where } T(b,0) = 6N_b^{2}+54N_b-12b .
\]
\end{theorem}

The load moves the right block $(111)^{1+b}$ into a new inner counter as the
value $2^{b}-1$, and takes one off the outer counter. The ticks then count the
inner counter down to zero, which builds a right block of $2^{b}-1$ units.
These are the lemmas \coq{Rst2}, \coq{Rst1}, \coq{Rst0} of \cite{mxdys} (called
Inc2, Inc1, Inc0 on the wiki \cite{wiki}). In their numbers the load costs
$12a + 3b + 12j_0 + 63 + 4m$ steps. Since $q$ grows by $3b+5$, the marker
$m = q \bmod 3$ goes $2 \to 1 \to 0 \to 2$ from one epoch to the next.

\section{Restart and halt}

When the outer counter is zero, the head runs to the left end of the tape. What
happens next depends on $q \bmod 3$.

\begin{theorem}[Restart; \coq{restart}]\label{thm:restart}
Let the outer counter be $L$ zero bits and let $q = 3a+1$. Then
\[
S_1(0^{L},\ 3a+1,\ b) \;\longrightarrow\; S_1\bigl(\beta^{*},\ 3b+2,\ 2^{b}-1\bigr)
\]
in exactly $12L + 12a + 3b + 97 + T(b,0)$ steps, where the new outer counter
$\beta^{*}$ has $L+a+2$ bits: $L+1$ ones, then a zero, then $a$ ones. Its value
is $2^{L+a+2} - 2^{a} - 1$.
\end{theorem}

\begin{theorem}[Halt; \coq{last\_move\_cfg}, \coq{halt\_ones}]\label{thm:halt}
Let the outer counter be $L$ zero bits and let $q = 3a$. From $S_1(0^{L}, 3a, b)$
the machine halts after exactly
\[
12L + 12a + 60 \quad\text{steps,}
\]
and the final tape holds exactly $L + a + 3b + 16$ ones.
\end{theorem}

The count of ones can be read off the final tape: $6$ in the left end
$1110111$, $L+4$ in what was the outer counter, $a+1$ in the middle block, $2$
written just before the halt, and $3b+3$ in the right block, which is not
touched.

\section{The whole run}

\begin{proposition}[\coq{init}]\label{prop:init}
From the blank tape the machine reaches $S_1(111,\ 20,\ 63)$ after exactly
$27{,}965$ steps. In the numbers of mxdys this is $S_1(3,7,2,6,0,63)$.
\end{proposition}

By Theorem~\ref{thm:epoch} the first epoch then takes
$332 + T(63,0)$ steps. It ends at step
\[
510{,}423{,}550{,}381{,}407{,}695{,}582{,}443{,}536{,}695{,}552{,}928{,}613
\approx 5.1\times 10^{38}
\]
(\coq{first\_epoch}). The second epoch has $2^{2^{63}-1}-1$ ticks, so no later
step number can be written out.

\paragraph{Many epochs.} Between restarts only two numbers change, $q$ and $b$,
together with the outer counter. One epoch is the map
$(q,b)\mapsto(q+3b+5,\ 2^{b}-1)$ with the cost of Theorem~\ref{thm:epoch}. The
Coq file defines two functions that add up the costs of a long run of epochs by
recursion on the shape of the outer counter, in the same way as the proof of
Theorem~\ref{thm:count}: \coq{runF}~$k$ clears $k$ low bits that are all $1$
($2^{k}-1$ epochs), and \coq{runTop}~$n$~$z$ clears $n$ ones that sit above $z$
zeros. The theorems \coq{runF\_ok} and \coq{runTop\_ok} say that the machine
follows these functions exactly.

\paragraph{Restart or halt.} Two simple facts decide the whole run
(\coq{runF\_inv}, \coq{runTop\_inv}): $b$ is always odd, and therefore each
epoch adds $2$ to $q$ modulo $6$. Together with $N_k \equiv 0$ or $1 \pmod 3$
for $k$ even or odd, this gives the three passes found by mxdys:
\begin{center}
\begin{tabular}{clll}
\toprule
Pass & Outer counter at the start & Epochs mod 3 & $q$ at the end\\
\midrule
1 & $111$ (value 7) & $1$ & $6t+4 = 3a_1+1$, $a_1$ odd: restart\\
2 & $2^{a_1+5}-2^{a_1}-1$ & $1$ & $6t+7 = 3a_2+1$, $a_2$ even: restart\\
3 & $2^{a_1+a_2+7}-2^{a_2}-1$ & $2$ & $6t+9 = 3a_3$: halt\\
\bottomrule
\end{tabular}
\end{center}

\begin{theorem}[Exact halting time; \coq{bb6\_halts\_exactly}]\label{thm:main}
Let $H$ be the sum of: $27{,}965$; the costs of the $7$ epochs of pass~1; the
first restart; the costs of all epochs of pass~2; the second restart; the costs
of all epochs of pass~3; and the final move $12L_3+12a_3+60$, where
$L_3 = a_1+a_2+7$. Then the machine, started on the blank tape, is halted after
$H$ steps and is not halted after any smaller number of steps.
\end{theorem}

In Coq, $H$ is the closed term \coq{halting\_time}, built from \coq{runF},
\coq{runTop}, $T$ and the formulas of Theorems
\ref{thm:epoch}--\ref{thm:halt}. Coq never evaluates it.

\begin{theorem}[Final ones; \coq{bb6\_final\_ones}]\label{thm:ones}
The final tape holds exactly $L_3 + a_3 + 3b_3 + 16$ ones, where $b_3$ is the
value of $b$ at the last shape.
\end{theorem}

Since $b_3$ is far larger than $L_3$ and $a_3$, the number of ones is very
nearly $3b_3$. This agrees in size with the estimates on the wiki \cite{wiki}.

\paragraph{Runtime against ones.} Let $b'$ be the value of $b$ at the start of
the last epoch, so that $b_3 = 2^{b'}-1$. By Theorem~\ref{thm:epoch} the ticks
of that one epoch take $6b_3^{2}+54b_3-12b'$ steps. So, writing $\sigma$ for the
number of ones,
\[
H \;=\; 6b_3^{2} + 54b_3 + R, \qquad \sigma \;=\; 3b_3 + r, \qquad r = L_3+a_3+16,
\]
where $R$ collects every other step of the run. Both $R$ and $r$ are tiny
beside $b_3$: each earlier epoch is smaller than the next by a whole
exponential, so all of them together cost about $(\log b_3)^{2}$ steps, and
$L_3$ and $a_3$ are about $\log b_3$ or less. Hence
\[
H \;\approx\; \tfrac{2}{3}\,\sigma^{2} .
\]
The two displayed identities follow from the theorems above. The size of $R$
and $r$ is a pen-and-paper estimate and is not stated in Coq. As a small test,
a starting shape whose run ends with $b=127$ gives $H/\sigma^{2}=0.61$ and
$R = 1044$. A running time of the order of $\sigma^{2}$ is what one expects of
a machine that sweeps back and forth over its tape.

\section{Transcript words}\label{sec:words}

The moves above were first seen in the transcript. Write $\Omega=\w{479}$ and
$\Lambda=\w{5b9}$. The transcript of each move is one word with a few counts:

\begin{center}
\renewcommand{\arraystretch}{1.25}
\begin{tabular}{@{}l>{\raggedright\arraybackslash}p{0.74\textwidth}@{}}
\toprule
Move & Word\\
\midrule
Tick, $j\ge1$ &
$\w9\,\Omega^{3j-1}\,\w{478602}\,\Lambda^{j+1}\,\w{4786025a}\,\w1^{3b+3}\,\w{02}\,\Omega^{3b+5}\,\w{47879}\,\Omega^{2}\,\w{47}$\\
Tick, $j=0$ &
$\w{8602}\,\Lambda\,\w{4786025a}\,\w1^{3b+3}\,\w{02}\,\Omega^{3b+5}\,\w{47879}\,\Omega^{2}\,\w{47}$\\
Load &
$\w9\,\Omega^{q-1}\,\w{47878602}\,\Omega^{3j_0+3}\,\w{478602}\,\Lambda^{j_0+1}\,\w{5a}\,\w1^{q+3}\,\w{02}\,\Lambda^{b+1}\,\Phi$\\
Restart &
$\w9\,\Omega^{q-1}\,\w{47878602}\,\Pi\,\Lambda^{a+1}\,\w{5b8602}\,\Lambda^{b+1}\,\Phi$\\
Halt &
$\w9\,\Omega^{q-1}\,\w{47878602}\,\Pi\,\Lambda^{a+1}\,\w{5a03}$\\
\bottomrule
\end{tabular}
\end{center}

\noindent Here
\begin{align*}
\Phi &= \w{4786025a02}\,\Omega^{2}\,\w{47879}\,\Omega^{2}\,\w{47},\\
\Pi &= \Omega^{3L+5}\,\w{47879}\,\Omega\,\w{4786025b8602}\,\Lambda^{L+2}\,\w{478602}.
\end{align*}
The word $\Phi$ closes the load and the restart, and $\Pi$
is the turn at the left end of the tape, and for $q=0$ the opening
$\w9\,\Omega^{q-1}\,\w{47878602}$ is replaced by $\w{878602}$. The lengths of
these words are the step counts of the theorems. The halt word ends in $\w3$,
the only occurrence of the halting pair in the whole run.

The words were checked against direct simulation on $228$ cases with small
parameters. The Coq proofs follow exactly these words, one lemma for each
repeated word and a computation for each fixed piece, but the digit strings are
not yet stated as theorems.

\begin{remark}
In the first $30{,}000{,}000$ steps from the blank tape, the digit $\w{a}$ cuts
the transcript into $2{,}305$ pieces. Of these, $2{,}299$ are exactly one tick word each (cut at a different point
of the word), with the counts given above. The other six lie in the first
$800$ steps and around the first load. This is how the tick was found, before
the proof of mxdys was read.
\end{remark}

\section{The Coq proof}

The development is the file \texttt{bb6/BB6Steps.v} in the repository (about
890 lines, Coq 8.18). \texttt{Print Assumptions} reports that
\coq{bb6\_halts\_exactly}, \coq{bb6\_final\_ones}, \coq{first\_epoch} and
\coq{T\_closed} are closed under the global context: no axioms are used.

\paragraph{Model.} As in \cite{bb5}, a configuration is a state, the list of
cells to the left of the head (nearest first), the cell under the head, and the
list of cells to the right. An empty list stands for blanks. \coq{steps n c}
applies the transition table $n$ times.

\paragraph{Structure.} Three lemmas handle the repeated words $\w1$, $\w{479}$
and $\w{5b9}$ for any number of repeats, by induction. Each move is then a
chain of these lemmas and of short fixed computations, with the step counts
added along the chain. Counters are lists of bits, so that a countdown is an
induction on the list. The numbers $a_1,a_2,a_3$ are tower-sized, but they only
ever appear as variables.

\paragraph{The model is the right machine.} The start of the run is checked by
computation inside Coq: $821$ steps from the blank tape give a shape $S_0$, and
Theorem~\ref{thm:count} then gives Proposition~\ref{prop:init}.

\paragraph{Independent test.} A Python script in the repository,
\texttt{check\_steps.py}, starts the machine from $132$ small shapes $S_1$ that
halt, and predicts the halting step and the final ones using only the formulas
of this paper and the rules of mxdys. All $132$ agree with plain simulation.

\paragraph{What is not in Coq.} The transcript words of
Section~\ref{sec:words}. Also, this proof and that of mxdys use different tape
models and are not linked to each other; they agree on the shapes and on the
three passes.

\section{Discussion}

\paragraph{What kind of result this is.} Nothing here changes what is known
about the behaviour of the machine. The result is exact bookkeeping on top of a
known structure: the proof of halting said that the machine stops, and this
paper says after how many steps, as a definite expression. We do not know of an
earlier exact expression for the halting time of this machine, but one may
exist in unpublished notes or code of the bbchallenge community.

\paragraph{The step counts are simple.} Every move costs a linear number of
steps in its parameters, and a countdown costs a quadratic in $2^{b}$. The
enormous running time comes only from the rule $b \mapsto 2^{b}-1$, applied
once per epoch.

\paragraph{Transcripts as a way to find rules.} In \cite{bb5} the transcript of
the BB(5) champion was described by a grammar of repeated words with growing
counts. The same reading of the BB(6) champion's transcript gave the tick, its
two counts and its cost before any tape analysis was used. For the later moves
this is not possible: the first epoch alone takes about $5\times10^{38}$ steps,
so no simulation reaches a restart. There the rules of mxdys were needed to
know where to look, and small starting shapes were used to measure the costs.

\paragraph{Future work.} (1) State the transcript words in Coq. (2) Derive
small facts about the halting time, such as its remainder modulo a small
number. (3) Apply the same counting to the other machines covered by
\cite{mxdys}.

\section*{Acknowledgements}

The tape shapes and rules used throughout are due to mxdys. A question from
Shawn Ligocki led to the paragraph on runtime against ones. Parts of the
analysis and of the Coq development were carried out with the help of an AI
assistant (Claude, Anthropic); the Coq proofs are checked by the Coq kernel.

\begin{thebibliography}{9}

\bibitem{mxdys}
mxdys. \emph{busycoq}, file \texttt{verify/SOBCv5.v}, module \texttt{TM62}.
\url{https://github.com/ccz181078/busycoq}, commit \texttt{3f302b8}.

\bibitem{wiki}
BusyBeaverWiki. \emph{1RB1RA\_1RC1RZ\_1LD0RF\_1RA0LE\_0LD1RC\_1RA0RE}.
\url{https://wiki.bbchallenge.org/wiki/1RB1RA_1RC1RZ_1LD0RF_1RA0LE_0LD1RC_1RA0RE},
retrieved November 2026.

\bibitem{bb5}
V. Vaidyanathan. \emph{A proved transcript grammar for the BB(5) champion}.
November 2026.
\url{https://github.com/VishnuVaidyanathan/BB5-decimal-macrogrammar}.

\bibitem{bbchallenge}
The bbchallenge Collaboration. \emph{Determination of the fifth Busy Beaver
value}. arXiv:2509.12337, 2025.

\bibitem{rado}
T. Rad\'o. On non-computable functions. \emph{Bell System Technical Journal},
41(3):877--884, 1962.

\end{thebibliography}

\end{document}
